\input{libraries/report-preamble.tex}
\title{Euclidean GCD on Unsigned 64-Bit Integers}
\begin{document}
\maketitle

\section{Task and interface}
\texttt{LeanExe.Lib.NumberTheory.gcd} computes the greatest common divisor
of two \texttt{UInt64} values using Euclid's remainder algorithm.  It accepts
every pair of input words.  The intended result is $\gcd(a,b)$, represented
as a word, with $\gcd(0,0)=0$.

LeanExe compiles the Lean function and its callers to WebAssembly (WASM).
The examples run that WASM in Wasmtime.

\section{Verification status}
Lean checks the source types.  Execution tests pass for selected inputs:
WASM and native Lean results agree with an arbitrary-precision integer
reference, and the native test also checks Lean's \texttt{Nat.gcd}.
Formal correctness and termination proofs for this source function,
its generated WASM, and its clients remain deferred.

\section{Algorithm and source}
Each iteration replaces $(x,y)$ with $(y,x\bmod y)$ while $y>0$.
The recurrence preserves the common divisors, and the nonzero second operand
decreases.  When it reaches zero, the first operand is the result.  NIST's
algorithm entry gives this recurrence~\cite{euclid}.  The zero-input convention
extends its positive-input description.

\lstinputlisting{LeanExe/Lib/NumberTheory/Gcd/Basic.lean}

The listing comes from \texttt{Basic.lean}.
The loop performs one unsigned remainder per iteration.  For $(48,18)$,
the successive pairs are $(48,18)$, $(18,12)$, $(12,6)$, and $(6,0)$.
This gives result $6$ after three iterations.

\section{Clients and execution}
The fraction client calls the component and divides both inputs by the returned
divisor when the denominator is positive.  Its result for $48/18$
is $8/3$.  A zero numerator yields $0/1$.

The polynomial-ratio client calls Horner evaluation with overflow checks for its
numerator and denominator, then calls fraction reduction.  At $x=2$,
$(1+3x+2x^2)/(2+2x)=15/6=5/2$.

The commands run from the repository root after following the
\href{https://github.com/jsmorph/leanexe/blob/lib1/DEVELOPING.md}{setup guide}.

\begin{lstlisting}
tools/seminum gcd 48 18
6
tools/seminum fraction 48 18
8/3
tools/seminum ratio 2 1,3,2 2,2
5/2
\end{lstlisting}

The command-line parser accepts unsigned decimal words through $2^{64}-1$.
Invalid syntax and a zero fraction denominator produce status 2 and a message
on standard error.  Successful commands print their result on standard output.

\section{Tests}
From the repository root, \texttt{node test/seminum.js} runs the comparisons
described above.  Cases include zero inputs, maximum words, powers of two,
consecutive large Fibonacci numbers, and generated full-width pairs.
Client cases check fraction reduction, polynomial ratios, and rejection of
a zero denominator.  CLI tests check results and invalid-input diagnostics.
The cases are defined in \texttt{test/seminum.js} and the native reference
driver \texttt{test/SeminumNative.lean}.

\texttt{tools/seminum inspect} lists the component declarations found in the
compiler annotations, with their catalog descriptions and proof status.

\begin{thebibliography}{9}
\bibitem{euclid} Paul E. Black, ``Euclid's algorithm,'' NIST Dictionary of
Algorithms and Data Structures, 2015.
\url{https://xlinux.nist.gov/dads/HTML/euclidalgo.html}.
The recurrence is the algorithm used here.  The component specifies zero inputs
and a fixed-width representation in addition to that reference.
\end{thebibliography}
\end{document}
